\section{Related Work}
\label{sec:related}

\para{Internal representation probing.}
A growing body of work shows that LLMs encode truthfulness in their hidden states.
\citet{azaria2023internal} train classifiers on activation vectors to detect when an LLM ``knows'' it is producing false output, achieving over 90\% accuracy across six factual domains.
\citet{marks2023geometry} discover that truthfulness is represented as a linear direction in LLaMA's activation space, suggesting a consistent internal geometry of truth.
\citet{burger2024truth} extend this finding with a two-dimensional truth subspace that generalizes across topics, models, and languages, reaching 93.8\% accuracy on real-world lying scenarios.
These results establish that truth and falsehood are internally distinguishable, but they require white-box access to model activations---unavailable for closed-source APIs.
Our work asks whether this internal distinction surfaces in the text itself.

\para{Behavioral lie detection.}
\citet{pacchiardi2024catch} propose a black-box approach: after a suspected lie, they ask 48 follow-up ``elicitation questions'' and classify the response pattern, achieving AUC of 1.0 in-distribution and 0.89 cross-architecture.
This demonstrates that lying creates a detectable behavioral fingerprint extending beyond the initial response.
However, this method requires multi-turn interaction, which is not always feasible.
We complement this line of work by analyzing whether a \emph{single} response contains sufficient stylistic signal to detect deception.

\para{Strategic deception in LLMs.}
Recent studies show that LLMs can engage in deception strategically, not only when explicitly instructed.
\citet{scheurer2023deception} demonstrate that LLMs deceive in agentic scenarios (such as insider trading simulations) without explicit instruction, while \citet{ogara2023hoodwinked} test LLM deception in interactive social deduction games.
\citet{wang2025thinking} study how chain-of-thought reasoning models formulate deceptive strategies with detectable patterns in their reasoning traces.
These findings motivate the study of deceptive language style as a potential detection mechanism.

\para{Stylometry and machine-generated text.}
\citet{schuster2019limitations} directly address whether stylometric features can detect deceptive intent in machine-generated text.
Using GPT-2-era models, they find that while stylometry can distinguish human from machine text (provenance detection), it cannot distinguish deceptive from legitimate machine text (intent detection).
They conclude that language models generate stylistically consistent text regardless of the content's veracity.
Our work revisits this conclusion with modern instruction-tuned models and finds the opposite result: RLHF-trained models \emph{do} exhibit measurable stylistic shifts when instructed to lie.

\para{Deception cues in language.}
\citet{hazra2024tell} identify four linguistic cues of deception---ambiguity, overconfidence, half-truths, and entailment violations---using a concept-bottleneck model applied to human deception data from the ``To Tell The Truth'' game show.
\citet{jones2024lies} examine LLMs' capacity for persuasion and deception through distributional language statistics.
Our feature set draws on established deception cue categories from psychology \citep{newman2003lying}, adapted for LLM-generated text with features targeting hedging, certainty, and structural patterns.
